\documentclass{article}

\usepackage[sglblindworkshop]{template/neurips_2026}
\workshoptitle{AI for Peace Workshop}
\makeatletter
\renewcommand{\@noticestring}{Prepared for the NeurIPS 2026 AI for Peace Workshop.}
\makeatother

\usepackage[utf8]{inputenc}
\usepackage[T1]{fontenc}
\usepackage{hyperref}
\usepackage{microtype}
\hypersetup{hidelinks}

\title{Matched Sources Expose Distinct Admissibility Failures in Small-Model Appeal Review}
\author{Rafael Gardos\\\texttt{gardos.rafael@gmail.com}}

\begin{document}
\nolinenumbers

\maketitle

\begin{abstract}
Plausible evidence can still be inadmissible under a service policy. We pair 48 synthetic humanitarian appeals from unlisted sources with 24 controls from accepted sources. Qwen and Gemma accept every record in both classes, giving 50.0 percent balanced source discrimination. Phi rejects every unlisted record and accepts half of the controls, reaching 75.0 percent. Mistral reaches 85.4 percent by accepting 77.1 percent of valid records and rejecting 93.8 percent of unlisted records. Under the printed policy, exact invalid-source accuracy is 0 of 384. An alternative information-request target raises Mistral's independent balanced exact score from 38.5 to 85.4 percent and Phi's from 25.0 to 62.5 percent. Qwen and Gemma remain constant acceptors under both conventions. Earlier-rationale exposure raises Mistral's false eligibility by 35.42 points, with a 12-base-request interval of [29.17, 41.67]. Four pinned four-bit models produce 1,152 matched reviews over two prompt orders. The fictional cases support no operational aid decision. For peace-oriented review, we define the three-step Matched Admissibility Audit. It pairs accepted and unlisted sources. It checks exact dispositions under both scoring conventions, then repeats the comparison after prompt-order changes.
\end{abstract}

\end{document}
